%% ma: L10-B14-0049
%% lop: 10
%% chuong: 5
%% bai: 14
%% dang: 10.14.8
%% loai: TL
%% muc: TH
%% dap_an: "-"
%% nguon: "chuyen-de-cac-so-dac-trung-cua-mau-so-lieu-khong-ghep-nhom-toan-10-chuong-trinh-moi.pdf (GV. Phan Nhật Linh) (P05), Bài 03 Dạng 2 Bài tập tự luận 9"
%% kien_thuc: "Lớp 10 Bài 13 (số trung bình), Bài 14 (phương sai, độ lệch chuẩn)"
%% trang_thai: chinh-thuc
%% ghi_chu: "Tự duyệt theo yêu cầu giáo viên 22/10/2026."
\begin{ex}
Điểm kiểm tra chất lượng đầu năm môn Toán (thang điểm 10) của một lớp có $40$ học sinh cho bởi bảng sau:
\begin{center}
\begin{tabular}{|c|c|c|c|c|c|c|c|c|}\hline
Điểm & $3$ & $4$ & $5$ & $6$ & $7$ & $8$ & $9$ & $10$\\\hline
Tần số & $2$ & $3$ & $9$ & $4$ & $8$ & $10$ & $3$ & $1$\\\hline
\end{tabular}
\end{center}
Tính độ lệch chuẩn của mẫu số liệu (làm tròn đến hàng phần trăm).
\loigiai{$\overline x=\dfrac{260}{40}=6{,}5$.

$s^2=\dfrac{2(3{,}5)^2+3(2{,}5)^2+9(1{,}5)^2+4(0{,}5)^2+8(0{,}5)^2+10(1{,}5)^2+3(2{,}5)^2+(3{,}5)^2}{40}=\dfrac{120}{40}=3$.

$s=\sqrt3\approx1{,}73$.}
\end{ex}
